\section{The protected background transpose and its critical tangent currents}
\label{sec:r13-protected-background-trace}

The positive background in the remaining threefold Dirac
problem has degree one.  On a protected finite moment domain
its exact action on pulled-back smooth tests is a positive
closed real supercurrent of bidimension two.  Its trace has
quadratic growth.  This supplies a uniform critical bound,
without assuming a local potential for the background or
discarding any part of a moment fiber.

\subsection{An exact transpose on a complete proper moment domain}

Let \(U\) be an original ordinary domain of pure dimension
three, and let \(t:U\to\Omega\subset\mathbf R^N\), \(N\ge2\), be one
finite actual moment map on a GP-protected domain.  Thus
\(\Omega\) is open, and
\[
       t^{-1}(L)\text{ is compact in }U
                         \quad(L\Subset\Omega\text{ compact}).
\tag{MT1}
\]
The inverse image in MT1 is the entire ordinary inverse
image.  It is not an old-source inverse image.  Let \(S\)
be any original symmetric middle-positive separately
\(d'\)- and \(d''\)-closed current of degree \((1,1)\)
on \(U\).  Define, for every compactly supported smooth
real \((2,2)\)-superform \(\eta\) on \(\Omega\),
\[
                       \langle T,\eta\rangle
                              =\langle S,t^*\eta\rangle.
\tag{MT2}
\]

\begin{proposition}[Middle positivity and both closedness equations]
\label{prop:r13-protected-background-transpose}
Formula MT2 defines a symmetric middle-positive real
supercurrent of bidimension \((2,2)\), hence bidegree
\((N-2,N-2)\), with
\[
                       d'T=d''T=0.
\tag{MT3}
\]
For \(\beta=\tfrac12D|x|^2\), its positive trace measure
is
\[
             \mu_T=\tfrac12T\wedge\beta^2,
 \qquad
 \mu_T(A)=\tfrac12
       \bigl(S\wedge t^*\beta^2\bigr)(t^{-1}(A))
\tag{MT4}
\]
for relatively compact Borel subsets \(A\) of \(\Omega\).
Every coefficient measure of \(T\) has total variation
bounded by a fixed dimensional multiple of \(\mu_T\).
\end{proposition}

\begin{proof}
MT1 makes \(t^*\eta\) an original compact smooth CLD
test.  On each fixed compact support the finite pullback
is continuous in the usual smooth test topology, proving
that MT2 is a current.  Column exchange commutes with
this pullback, giving symmetry.  If \(\alpha\) is a
smooth real \((2,0)\)-form, the standard signed form
\(\alpha\wedge J\alpha\) is middle-positive and so is
its pullback.  The assumed original middle positivity
of \(S\) therefore gives exactly middle positivity of
\(T\).  No equivalence with weak or strong positivity
in the ambient dimension \(N\) has been used.

Pullback commutes with both first derivatives.  Compact
support remains inside \(U\) by MT1, so applying either
closedness equation of \(S\) to a pulled-back test proves
MT3.  In particular no source cutoff has been inserted:
such a cutoff could have introduced an additional boundary
derivative, whereas MT2 uses the whole protected inverse
image.

The scalar functional obtained by testing against
\(\varphi\beta^2/2\) is a positive Radon measure.
The equality for smooth \(\varphi\) in MT2 identifies
it with the proper measure transpose in MT4, and hence
gives that equality for Borel sets.  In a constant
\((2,0)\) basis the cocoefficient measures form a
positive-semidefinite matrix.  Its diagonal sum is the
trace, and polarization bounds each off-diagonal
variation by one half the sum of the two corresponding
diagonal measures.  This gives the asserted coefficient
control, with the fixed orientation and basis constants.
\end{proof}

This construction is an exact finite-presentation test
transpose.  It does not assert a higher-degree analytic
action, a pushforward formula for an arbitrary cap carrier,
or that tests in this one target replace all original tests.
The current \(S\) still has its original complete carrier.
In particular MT4 includes background mass at every point
of the relevant moment fiber, regardless of its old rank.

\subsection{The exact quadratic trace estimate}

Translate \(t(y)\) to zero and fix
\(\overline{B_{r_0}(0)}\subset\Omega\).  Let
\(M_T(r)=\mu_T(B_r(0))\).  The following is the
bidimension-two case of the real Lelong monotonicity
calculation; the normalization is recorded explicitly.

\begin{theorem}[Quadratic growth and finite critical density]
\label{thm:r13-protected-quadratic-trace-growth}
The function
\[
                         r\longmapsto\frac{M_T(r)}{r^2}
\tag{MT5}
\]
is nondecreasing for \(0<r<r_0\).  Consequently
\[
 \begin{split}
 M_T(r)&\le \frac{M_T(r_0)}{r_0^2}\,r^2,\qquad 0<r<r_0,\\
 \nu(T,0)&=\lim_{r\downarrow0}\frac{M_T(r)}{r^2}
                                      \quad\text{exists and is finite}.
 \end{split}
\tag{MT6}
\]
No assertion that \(\nu(T,0)=0\) is part of the theorem.
In particular \(\mu_T(\{0\})=0\), while the density
in MT6 can be positive.
\end{theorem}

\begin{proof}
First suppose \(T\) is smooth.  On an annulus write
\(r=|x|\) and use
\[
 D r=\frac{\beta}{r}
                -\frac{d'r\wedge d''r}{r}.
\tag{MT7}
\]
On the sphere of radius \(r\), the first differential
\(d'r\) pulls back to zero.  Stokes, using \(d'T=0\),
therefore gives for \(0<a<b<r_0\)
\[
 \frac{M_T(b)}{b^2}-\frac{M_T(a)}{a^2}
             =\frac12\int_{a<|x|<b}T\wedge(D|x|)^2\ge0.
\tag{MT8}
\]
To check the constant, apply Stokes first to
\(T\wedge\beta^2\), and note on a sphere that
\(d''(|x|^2/2)=r\,d''r\) and
\(\beta=r\,Dr\) after tangential pullback in
the Stokes boundary expression.  The resulting boundary
identity is
\[
 \frac{M_T(r)}{r^2}
       =\frac12\int_{|x|=r}T\wedge d''r\wedge Dr.
\]
Subtract this identity at \(a\) from that at \(b\),
and apply Stokes on the annulus, where \(r\) is smooth.
This proves MT8 without introducing any product of
\(D|x|\) at the origin.
The integrand on the annulus is positive: \(Dr\)
is semipositive and its square is strongly positive,
so it may be paired with the middle-positive \(T\).

For a general coefficient-Radon current convolve its
coefficients by a nonnegative compact smoothing kernel,
inside a slightly smaller ball.  The resulting currents
are smooth, middle-positive and separately closed.
Their trace measures converge weakly to \(\mu_T\).
At radii whose spheres have zero trace mass, MT8
passes to the limit.  There are at most countably
many exceptional radii.  Approximation of arbitrary
open balls from below by nonexceptional radii gives
the monotonicity in MT5 for all \(a<b\).  The finite
measure on the fixed larger compact ball gives MT6;
letting \(r\downarrow0\) also gives the absence of
an atom of the trace itself.
\end{proof}

The printed statement of Lagerberg's Proposition 3.4
omits the word closed, but its Stokes proof and the
following corollary use closedness.  MT3 is imposed
and proved here.  Middle positivity is also essential
for coefficient control: Lagerberg's Remark 3.3
exhibits a weakly positive nonzero tensor with zero
trace.  No such weaker cone is substituted in MT4.

The exponent two is sharp.  A positive integration
current of a complete real two-plane has
\(M_T(r)=c r^2\), with \(c>0\).  Thus MT6
supplies the critical bound required for a first-jet
replacement in CR; it does not supply CR5's little-o
bound for the leading homogeneous term.

\subsection{Compactness and conicity of critical tangent currents}

For \(\delta>0\), let \(a_\delta(x)=x/\delta\)
and define on \(a_\delta(\Omega)\)
\[
                  T_\delta=\delta^2(a_\delta)_*T.
\tag{MT9}
\]
The pushforward here is only the ordinary real dilation
of a real supercurrent.  Equivalently, each coefficient
measure transforms as
\[
                    (\mu_{IJ})_\delta(A)
                                =\delta^{-2}\mu_{IJ}(\delta A).
\tag{MT10}
\]
The factor in MT9 accounts for both columns of a
\((2,2)\) test; it is not a change of the original
non-Archimedean absolute value.

\begin{corollary}[Critical tangents are positive closed cones]
\label{cor:r13-critical-tangent-current-compactness}
For every sequence \(\delta_j\downarrow0\), a subsequence
of \(T_{\delta_j}\) converges on all compact smooth
real tests to a global middle-positive separately
closed supercurrent \(T_0\) of bidimension two.
Every such limit satisfies
\[
       \mu_{T_0}(B_R(0))=\nu(T,0)R^2\quad(R>0),
       \qquad \lambda^2(a_\lambda)_*T_0=T_0
                                      \quad(\lambda>0).
\tag{MT11}
\]
When the original complete moment-image germ is contained
in a closed cone \(C\), each limit is supported in \(C\).
The tangent limit need not be unique or polyhedral.
\end{corollary}

\begin{proof}
For any fixed \(R>0\), MT6 gives
\[
 \mu_{T_\delta}(B_R)
       =\delta^{-2}M_T(\delta R)
       \le R^2 M_T(r_0)/r_0^2
\tag{MT12}
\]
for small \(\delta\).  MT4 bounds every coefficient
measure on each compact ball.  Weak compactness of
Radon measures and a diagonal selection on increasing
balls give a subsequential current limit.  Positivity,
symmetry and both closedness equations pass to smooth
tests.  Cone support is preserved under dilation and
under this limit.

The left side of MT12 converges numerically to
\(\nu(T,0)R^2\).  The open- and closed-ball
inequalities for weak measure convergence, followed
by \(R-\epsilon\uparrow R\) and
\(R+\epsilon\downarrow R\), show that the limit
measure has exactly that ball mass.  In particular
it charges no sphere and no atom at the origin.

Apply the annular identity MT8 to \(T_0\), by the
same local regularization argument.  Its constant
normalized ball mass gives
\[
                T_0\wedge(D|x|)^2=0
                                 \quad\text{on }\mathbf R^N\setminus\{0\}.
\tag{MT13}
\]
In an orthonormal frame whose first vector is radial,
\(D|x|\) is \(|x|^{-1}\) times the identity on the
orthogonal hyperplane and is zero on the radial vector.
The middle-positive cocoefficient matrix of \(T_0\)
is positive semidefinite on two-vectors.  MT13 makes
its whole tangential two-vector block zero; positivity
then makes every row and column meeting that block
zero.  Only two-vectors containing the radial direction
remain.  In current-form notation this is
\[
                    \iota'_E T_0=\iota''_E T_0=0,
                      \qquad E=\sum_i x_i\partial_{x_i}.
\tag{MT14}
\]
The origin has zero coefficient mass, so these identities
hold globally.  The one-column Cartan identity and
\(d'T_0=0\) now imply, for every distributional
coefficient \(\mu_{IJ}\) of bidegree \((N-2,N-2)\),
\[
             E\cdot\partial\mu_{IJ}+(N-2)\mu_{IJ}=0.
\tag{MT15}
\]
Equivalently
\(\int(E\varphi+2\varphi)\,d\mu_{IJ}=0\)
for every compact smooth scalar \(\varphi\).
Differentiating
\(\lambda^{-2}\int\varphi(x/\lambda)\,d\mu_{IJ}(x)\)
therefore gives zero.  This is exactly the dilation
identity in MT11, by MT10.
\end{proof}

The conicity conclusion constrains the possible critical
backgrounds, but it does not put them in the retained
analytic weak delta-background class \(\mathcal W_{\rm loc}\).
It also does not prove that their full paired action on
the leading finite-PS tensor vanishes.  Such a density
or annihilation assertion is a further theorem.

\paragraph{Primary sources and original inputs.}
The real trace and monotonicity formulas were checked
directly in A.~Lagerberg, \emph{Super currents and tropical
geometry}, arXiv:1008.2856v1, Definition 3.1,
Propositions 3.2 and 3.4, and Remark 3.3, pp.~19--20,
\url{https://arxiv.org/pdf/1008.2856}.  The closed-current
scope and exponent are also explicit in F.~Elkhadhra
and K.~Zahmoul, \emph{Lelong--Jensen formula,
Demailly--Lelong numbers and weighted degree of positive
supercurrents}, arXiv:1809.00992v2, Theorem 1 and
Corollary 2, p.~11,
\url{https://arxiv.org/pdf/1809.00992}.  MT2--MT4 use
the retained original GP compactness and middle-positive
coefficient-Radon interface.  The critical scaling and
conicity deduction are recorded above with their full
normalizations.
